% Source: tex/chapters/ch06/05_ソフトウェアエンジニアリング運用計画.tex
\subsubsection{ソフトウェア可用性、継続性、サービス水準}

可用性とは、必要なときにサービスが利用可能である度合いである。継続性とは、大きな障害や災害が発生してもサービスを維持または復旧できる能力である。サービス水準とは、可用性、性能、応答時間、復旧時間などについて合意された水準である。

運用担当技術者は、プロジェクト初期に非機能要求として定義された可用性と継続性を満たすため、適切なインフラを計画、設計、実装、試験する。可用性は測定され、記録され、予定外の停止は調査される。サービス報告では、合意されたサービス水準に対して、実際の運用がどうだったかを示す。

SLAは、運用サービスの義務を明確にするために有効である。SLAがなければ、利用者は「常に速く、絶対に止まらない」ことを期待し、提供者は「通常は動いている」と考えるかもしれない。合意された水準がなければ、運用の成否を客観的に判断しにくい。
